% Solutions/hints for ch12 exercises (assembled into Appendix C)
\begin{solution}{exr:ch12-tangent}
(a)~Let $\vect{t}$ be a point of contact of a tangent from the origin to the circle of radius $R$ around $\pos_{\mathrm{rel}}$. The radius $\pos_{\mathrm{rel}}-\vect{t}$ is perpendicular to the tangent, so the triangle with vertices $\vect{0}$, $\vect{t}$, $\pos_{\mathrm{rel}}$ has a right angle at $\vect{t}$, hypotenuse $d$ and the leg $R$ opposite to the angle $\theta$ at the origin; hence $\sin\theta=R/d$ and, by Pythagoras, $\norm{\vect{t}}=\sqrt{d^2-R^2}$. (b)~If $t\vel\in\disc{\pos_{\mathrm{rel}}}{R}$ for some $t>0$, then $(t/\lambda)(\lambda\vel)$ is the same point, so $\lambda\vel\in CC_{A|B}$ for every $\lambda>0$. For the velocity obstacle, $\lambda(\vel_B+\vel)=\vel_B+(\lambda\vel+(\lambda-1)\vel_B)$ lies in $\vel_B\oplus CC_{A|B}$ for all $\lambda>0$ and all $\vel\in CC_{A|B}$ only if $(\lambda-1)\vel_B$ is a difference of cone vectors for every $\lambda$, which fails for $\vel_B\ne\vect{0}$ (take $\lambda$ large and $\vel$ near a boundary ray). (c)~By (b) membership in $CC_{A|B}$ depends only on the direction of $\vel_A-\vel_B$. After truncation the collision must occur before $\ttc$, and the time to collision scales like $1/\norm{\vel_{\mathrm{rel}}}$, so a slow velocity in the cone may be allowed while a fast one in the same direction is forbidden.
\end{solution}
\begin{solution}{exr:ch12-ttc}
With $\pos=(6,2)$, $\vel=(3,0)$ and $R=2.5$: $\vel\cdot\vel=9$, $\pos\cdot\vel=18$, $\pos\cdot\pos-R^2=40-6.25=33.75$, so the quadratic is $9t^2-36t+33.75=0$ with $\Delta/4=324-303.75=20.25$ and $\sqrt{\Delta/4}=4.5$. The roots are $t=(18\mp4.5)/9$, i.e.\ $t_c=1.5$ and $t_2=2.5$: the discs overlap between $1.5$ and $2.5$. The closest approach is at $t^*=\pos\cdot\vel/(\vel\cdot\vel)=2$, where the relative position is $(0,2)$ and the distance $2<R$, consistent with the overlap interval being symmetric around $t^*$. For $\vel=(-3,0)$ the discriminant is the same but $\pos\cdot\vel=-18<0$: both roots are negative ($-1.5$ and $-2.5$), the ``collision'' lies in the past and \cref{thm:ch12-ttc} reports none. \code{time\_to\_collision((6,2),(3,0),2.5)} returns $1.5$ and \code{time\_to\_collision((6,2),(-3,0),2.5)} returns \code{inf}.
\end{solution}
\begin{solution}{exr:ch12-empty}
The apex $\vel_B=-s\,\pos_{\mathrm{rel}}/d$ lies on the axis of the wedge at distance $s$ from the origin, on the side away from $B$. The reachable disc $\disc{\vect{0}}{v_{\max}}$ is centred on the axis at distance $s>v_{\max}$ from the apex, so seen from the apex it subtends the half-angle $\arcsin(v_{\max}/s)$ (the same right triangle as in \cref{exr:ch12-tangent}). It lies entirely inside the wedge of half-angle $\arcsin(R/d)$ if and only if $v_{\max}/s\le R/d$, i.e.\ $s\ge v_{\max}d/R$. With $v_{\max}=1.5$, $d=7.07$, $R=1$ the threshold is $s=10.6\,\mathrm{m/s}$. Just above it every candidate collides; the fallback minimises $\norm{\vel-\vel^{\mathrm{pref}}}+\kappa/t_c$ and, for a large enough $\kappa$, picks the reachable velocity with the largest time to collision, which is the one closest to the boundary of the wedge: full speed perpendicular to the line of sight, slightly away from $B$. Truncation does not help because at $s\approx10.6$ the collision happens after about $(d-R)/(s+v)\approx0.5\,\mathrm{s}<\ttc$. For $s$ well below the threshold but $d$ small, the wedge is wide ($\theta\to90^\circ$ as $d\to R$), the feasible velocities are those pointing away from $B$ or nearly perpendicular to the line of sight, and the closer $B$ is, the harder the manoeuvre; a horizon of a few times $R/v$ makes the agent react while $\theta$ is still small.
\end{solution}
\begin{solution}{exr:ch12-coding}
Hints. Use \cref{alg:ch12-simulate} with synchronous updates and record positions and velocities at every step; the metrics are the minimum over time and pairs of the centre distance, the path length divided by the start--goal distance, the first time at which the goal is within a tolerance, and the count of steps at which the feasible set of \cref{alg:ch12-choose} was empty. Count sign changes of the velocity component perpendicular to the start--goal line. With the book's parameters ($\dt=0.1$, $\ttc=5$, $R=1$, speed $1$, $v_{\max}=1.5$) the circle scenario gives the numbers of \cref{tab:ch12-crossing}; with $\ttc=\infty$ agents facing a wall of static discs stop before reaching it, with $\ttc=1$ the agents of the circle scenario turn late and pairs collide from $n=4$ on, and $36$ directions produce visibly larger sign-change counts than $180$. Keep the scenario generator, the metrics and the plotting code separate from the avoidance rule, so that RVO and \orca can be plugged in as alternative implementations of \code{choose\_velocity} in \cref{ch:ch13}.
\end{solution}
